\title{DeepSeekMath: Pushing the Limits of Mathematical Reasoning in Open Language Models}

\begin{abstract}
Mathematical reasoning represents a substantial hurdle for language models because of its inherently complex and organized framework. Here, we present DeepSeekMath~7B, an extension of DeepSeek-Coder-Base-v1.5 7B, further pre-trained with 120B mathematics-focused tokens collected from Common Crawl, and supplemented with natural language and programming code. Without the assistance of external toolkits or ensemble methods, DeepSeekMath~7B attains a notable 51.7\% on the competitive MATH benchmark, narrowing the performance gap to models such as Gemini-Ultra and GPT-4. When evaluating self-consistency with 64 sampled responses, the model reaches 60.9\% accuracy on MATH. DeepSeekMath~’s strength in mathematical reasoning primarily results from two crucial elements: (1) extensive utilization of publicly available online data, processed through a sophisticated data selection workflow, and (2) the implementation of Group Relative Policy Optimization (GRPO)—a modified version of Proximal Policy Optimization (PPO)—which augments mathematical reasoning capability while making PPO more memory-efficient.
\end{abstract}

\section{Introduction}

Large language models (LLMs) have transformed approaches to mathematical reasoning within artificial intelligence, catalyzing notable progress across both quantitative reasoning benchmarks \citep{MATH} and benchmarks for geometric reasoning \citep{trinh2024solving}. Additionally, these systems have played a pivotal role in supporting human problem-solvers with advanced mathematical tasks \citep{tao}. Despite this progress, premier systems like GPT-4 \citep{gpt4} and Gemini-Ultra \citep{gemini} remain proprietary, while currently available open-source models lag considerably in terms of performance.

This paper presents DeepSeekMath, a specialized language model for mathematics that surpasses existing open-source models in mathematical reasoning and performs comparably to GPT-4 on scholarly benchmarks. Central to our approach is the development of the DeepSeekMath Corpus—a high-quality, large-scale pre-training dataset containing 120B math tokens. This corpus is derived from Common Crawl (CC) with the assistance of a fastText-based classifier \citep{joulin2016fasttext}. For its initial training, the classifier leverages samples from OpenWebMath \citep{openwebmath} as positive examples and contrasts them with a wide assortment of unrelated web pages used as negative examples. Once trained, the classifier identifies additional positive samples within the CC, which are subsequently refined through manual human annotation. This improved set is then utilized to retrain and further optimize the classifier. Our evaluation demonstrates the resulting corpus's high quality: the DeepSeekMath-Base 7B model achieves 64.2\% on the GSM8K benchmark \citep{gsm8k} and 36.2\% on the MATH competition dataset \citep{MATH}, exceeding the performance of Minerva 540B \citep{minerva}. The DeepSeekMath Corpus also possesses multilingual features, leading to enhanced outcomes on Chinese mathematical benchmarks \citep{wei2023cmath,agieval}. We propose that our methodology in constructing mathematical datasets offers a valuable foundation for future research, with ample opportunities for further advancements.

DeepSeekMath-Base is constructed by initializing with DeepSeek-Coder-Base-v1.5 7B \citep{deepseek-coder}, as our findings suggest that leveraging a code-focused pre-trained model offers superior performance over initializing from a standard large language model (LLM). Notably, the mathematical training phase not only augments the model’s proficiency in mathematical reasoning but also yields measurable improvements on general benchmarks such as MMLU \citep{mmlu} and BBH \citep{bbh}, reflecting a broader enhancement of the model’s reasoning faculties.

Subsequent to pre-training, we carry out mathematical instruction tuning on DeepSeekMath-Base, employing chain-of-thought \citep{cot}, program-of-thought \citep{pot,pal}, and tool-augmented reasoning datasets \citep{tora}. This results in DeepSeekMath-Instruct 7B, which surpasses all other models of equivalent size and rivals the performance of leading 70B open-source instruction-tuned models.

In addition, we propose Group Relative Policy Optimization (GRPO), an RL algorithm derived from Proximal Policy Optimization (PPO) \citep{schulman2017proximal}. GRPO omits the need for a critic model and instead calculates the baseline using group-level scores, thereby markedly decreasing the computational burden during training. Leveraging only a fraction of the English instruction tuning dataset, GRPO secures significant gains over the already strong DeepSeekMath-Instruct baseline—demonstrated by substantial improvements in both in-domain tasks (GSM8K: 82.9\% $\rightarrow$ 88.2\%, MATH: 46.8\% $\rightarrow$ 51.7\%) and out-of-domain problems (e.g., CMATH: 84.6\% $\rightarrow$ 88.8\%) in the RL phase. 

Moreover, we provide a unified framework for conceptualizing approaches such as Rejection Sampling Fine-Tuning (RFT) \citep{yuan2023scaling}, Direct Preference Optimization (DPO) \citep{dpo}, PPO, and GRPO, situating them as either direct or reduced forms of RL algorithms. Extensive empirical studies—comparing, for example, online versus offline optimization, outcome-based versus process-based supervision, and single-step versus iterative RL—are conducted to thoroughly analyze the critical mechanisms underpinning this paradigm. Finally, we elucidate the underlying reasons behind the observed performance gains from RL on instruction-tuned models and outline future directions for advancing more effective RL strategies within this unified framework.

\subsection{Contributions}
This work makes advancements in large-scale math pre-training and undertakes an in-depth examination of reinforcement learning methodologies.

\textbf{Math Pre-Training at Scale}

\begin{itemize}[topsep=0pt]
    \item We demonstrate that Common Crawl, a publicly available resource, contains substantial quantities of math-relevant information. Through the development of a robust data selection framework, we assemble the DeepSeekMath Corpus—a carefully curated dataset encompassing 120B tokens from web sources prioritized for mathematical content. This corpus is approximately seven times larger than the math datasets employed by Minerva \citep{minerva} and surpasses the recent OpenWebMath dataset \citep{openwebmath} by a factor of nine.

    \item The DeepSeekMath-Base 7B model, after pre-training, matches the performance level of Minerva 540B \citep{minerva}, suggesting that model size is not the sole determinant of proficiency in mathematical reasoning. Our results show that strong results can be achieved by compact models when trained on well-selected, high-quality data.

    \item We report results from a sequence of math training experiments, highlighting that initializing with code training prior to math tasks boosts the model's performance in solving mathematical questions—regardless of whether tool usage is permitted. These observations contribute evidence to the longstanding debate concerning the impact of code training on reasoning, specifically supporting its beneficial effect on mathematical reasoning.

    \item Contrary to common practice, supplementing training data with arXiv papers does not produce any meaningful gains across the suite of mathematical evaluation tasks used in this study.
\end{itemize}

\textbf{Exploration and Analysis of Reinforcement Learning}

\begin{itemize}[topsep=0pt]
    \item We present Group Relative Policy Optimization (GRPO), a reinforcement learning algorithm designed for efficiency and effectiveness. GRPO eliminates the need for a critic model, relying instead on group score-based baseline estimation. This modification yields considerable savings in training cost when contrasted with Proximal Policy Optimization (PPO).
    \item We show that GRPO markedly improves the outcomes of the instruction-tuned DeepSeekMath-Instruct model, using only instruction-tuning datasets. In addition, we detect notable improvements in the ability to generalize to out-of-domain data during the reinforcement learning phase.
    \item We establish a unified framework for analyzing diverse approaches, including RFT, DPO, PPO, and GRPO. Our empirical analysis encompasses various experimental setups—such as online versus offline learning, outcome-based versus process-based supervision, single-step versus iterative reinforcement learning, and more—in order to thoroughly explore the core components underlying this unified view.
    \item Drawing upon this comprehensive framework, we investigate the underlying factors that contribute to the success of reinforcement learning, and delineate a set of prospective research directions for achieving more effective reinforcement learning in large language models (LLMs).
\end{itemize}

\subsection{Summary of Evaluations and Metrics}

\begin{itemize}[topsep=0pt]
    \item \textbf{English and Chinese Mathematical Reasoning}: We systematically evaluate model performance on a broad array of English and Chinese benchmarks that span mathematical tasks from elementary to tertiary education levels. English benchmarks assessed include GSM8K \citep{gsm8k}, MATH \citep{MATH}, SAT \citep{llemma}, OCW Courses \citep{minerva}, and MMLU-STEM \citep{mmlu}. For Chinese, we use MGSM-zh \citep{mgsm}, CMATH \citep{wei2023cmath}, Gaokao-MathCloze \citep{agieval}, and Gaokao-MathQA \citep{agieval}. We measure the capability of models both to produce independent written solutions without auxiliary tools and to resolve problems by leveraging Python.
    
    DeepSeekMath-Base demonstrates competitive performance relative to the proprietary Minerva 540B \citep{minerva} on English datasets and consistently surpasses leading open-source base models (such as Mistral 7B \citep{mistral} and Llemma-34B \citep{llemma}), regardless of math-specific pre-training, frequently by a substantial margin. Its particularly strong performance on Chinese tasks can be attributed to our inclusion of diverse, high-quality mathematical pre-training data spanning multiple languages, in contrast to prior works \citep{minerva,llemma} that focus solely on English. Following further instruction tuning and reinforcement learning in mathematics, DeepSeekMath-Instruct and DeepSeekMath-RL set new open-source standards by surpassing the 50\% accuracy threshold on the challenging, competition-grade MATH benchmark.
\end{itemize}

\item \textbf{Formal Mathematics}: We assess DeepSeekMath-Base on the informal-to-formal theorem proving benchmark introduced by \citep{dsp_proof}, specifically utilizing the miniF2F dataset \citep{minif2f} and the Isabelle proof assistant \citep{isabelle}. The results indicate that DeepSeekMath-Base achieves robust performance in few-shot autoformalization settings.

\item \textbf{Natural Language Understanding, Reasoning, and Code}: In order to thoroughly evaluate models' proficiency in language comprehension, reasoning, and programming, we benchmark DeepSeekMath-Base against the Massive Multitask Language Understanding (MMLU) dataset \citep{mmlu}, encompassing 57 multiple-choice tasks from a wide array of disciplines; BIG-Bench Hard (BBH) \citep{bbh}, which contains 23 complex problems predominantly involving multi-step reasoning; as well as widely-adopted coding benchmarks HumanEval \citep{codex} and MBPP \citep{mbpp}. Math-focused pre-training is shown to improve both reasoning and language understanding capabilities.

\section{Math Pre-Training}

\subsection{Data Collection and Decontamination} This section details the construction process for the DeepSeekMath Corpus derived from Common Crawl. As illustrated in Figure \ref{fig:data_collect}, we employ an iterative pipeline designed to efficiently collect a large-scale mathematical text dataset from Common Crawl, beginning with a high-quality seed corpus of mathematical materials. Notably, this methodology can also be extended to other specialized areas, including code.

To begin, we select OpenWebMath \citep{openwebmath}, a reputable resource comprising quality mathematical text from the web, as our initial seed. Using this dataset, a fastText model \citep{joulin2016fasttext} is trained to identify additional web pages resembling the content of OpenWebMath. For training, we sample 500,000 positive instances from the seed and another 500,000 web pages from Common Crawl as negative examples. The open-source fastText toolkit\footnote{\url{https://fasttext.cc}} is utilized, with hyperparameters set as follows: a vector size of 256, learning rate of 0.1, maximum n-gram length of 3, a minimum of 3 occurrences per word, and 3 epochs of training. To further narrow down the massive Common Crawl, we perform deduplication both at the URL and near-content levels, bringing the dataset to 40B HTML documents. Applying the trained fastText model, we retrieve mathematical web content from this deduplicated corpus. To eliminate low-quality materials, the retrieved pages are ranked by their fastText scores, and only the highest-ranked subset is kept. We determine the optimal quantity of data by conducting pre-training evaluations on sets consisting of the top 40B, 80B, 120B, and 160B tokens, ultimately retaining the top 40B tokens after the first iteration.

Following the initial round of data collection, a significant number of mathematical web pages are still left uncollected. This is largely attributable to the limited diversity within the positive sample set used to train the fastText model. To address this limitation, we expand the initial seed dataset by sourcing additional mathematical web content, thereby enhancing the training set for the fastText model. The process begins with partitioning the entire Common Crawl dataset into mutually exclusive domains, each defined as a group of web pages sharing an identical base URL. For each domain, we determine the proportion of web pages retrieved during the first collection phase. Domains from which more than 10\% of pages were retrieved are marked as math-related (for instance, \url{mathoverflow.net}). 

We then proceed with manually curating URLs within these identified domains that are linked to mathematical content (e.g., \url{mathoverflow.net/questions}). Uncollected pages falling under these URLs are subsequently incorporated into the seed set. By broadening the diversity of positive samples in this manner, the next iteration yields an improved fastText model that is capable of identifying a larger quantity of mathematical data. Across four iterative rounds, this methodology results in the accumulation of 35.5 million mathematical web pages, amounting to 120 billion tokens. It is observed that, by the fourth iteration, roughly 98\% of the data overlaps with what was already obtained during the third iteration, leading us to discontinue further data collection at that stage.

In order to prevent overlap with benchmark datasets, we adopt the filtering approach of \cite{deepseek-coder}, eliminating web pages that include questions or answers from English mathematical evaluation benchmarks such as GSM8K \citep{gsm8k} and MATH \citep{MATH}, as well as Chinese datasets like CMATH \citep{wei2023cmath} and AGIEval \citep{agieval}. The filtration protocol entails excluding any text segment where a 10-gram string exactly matches a substring from any benchmark source. For those benchmark segments shorter than 10 grams but containing at least 3 grams, we apply an exact match removal to eliminate contaminated pages from the training set.

\subsection{Validating the Quality of the DeepSeekMath~Corpus}

We conduct pre-training studies to assess the quality of the DeepSeekMath~Corpus in comparison to several newly available math-specific corpora:

\begin{itemize}[topsep=0pt]
    \item \textbf{MathPile} \citep{mathpile}: A collection of sources yielding 8.9B tokens, including textbooks, Wikipedia, ProofWiki, CommonCrawl, StackExchange, and notably arXiv—which constitutes more than 85\% of its contents;
    \item \textbf{OpenWebMath} \citep{openwebmath}: A subset of CommonCrawl carefully filtered for mathematical relevance, comprising 13.6B tokens;
    \item \textbf{Proof-Pile-2} \citep{llemma}: This corpus brings together OpenWebMath, AlgebraicStack (10.3B tokens of mathematical code), and 28.0B tokens from arXiv papers. In evaluating Proof-Pile-2, we maintain the arXiv:Web:Code token ratio of 2:4:1 as specified in \cite{llemma}.
\end{itemize}

\subsubsection{Training Setting}
\label{sec:quality-policy}
We perform mathematical instruction tuning on a general-purpose language model comprising 1.3B parameters, architecturally consistent with the DeepSeek LLMs \citep{deepseek-llm}, herein referred to as DeepSeek-LLM 1.3B. Separate models are trained on individual mathematical datasets, each exposed to 150B tokens. The entirety of our experiments utilizes the efficient and lightweight HAI-LLM \citep{haillm} training infrastructure. Adhering to the training methodology established in DeepSeek LLMs, we adopt the AdamW optimizer \citep{adamW} with settings $\beta_1=0.9$, $\beta_2=0.95$, and $\mathrm{weight\_decay}=0.1$. The learning rate follows a multi-stage decay schedule: it ramps up to its maximum after 2,000 warmup steps, drops to 31.6\% of this maximum at 80\% completion, and further diminishes to 10.0\% by 90\% of the way through training. The peak learning rate is set at 5.3e-4, employing a batch size of 4M tokens and a 4K token context window.

\subsubsection{Evaluation Results}

\textbf{The DeepSeekMath~Corpus demonstrates superior quality, comprehensive multilingual coverage, and stands as the largest available resource.}

\begin{itemize}[topsep=0pt]
    \item \textbf{High-quality}:
    We assess model performance on 8 mathematical benchmarks utilizing few-shot chain-of-thought prompting \cite{cot}. Table \ref{tab:corpora_comparison} provides evidence that the DeepSeekMath~Corpus yields the best results among the compared datasets. Further, Figure \ref{fig:corpora_comparisons} illustrates that models trained with DeepSeekMath achieve higher accuracy than those trained on Proof-Pile-2 at the 50B token point (corresponding to a single epoch over Proof-Pile-2), underscoring the superior average data quality present in DeepSeekMath.
    \item \textbf{Multilingual}:
    The DeepSeekMath~Corpus contains a wide array of linguistic data, with English and Chinese forming the largest portions. As presented in Table \ref{tab:corpora_comparison}, leveraging DeepSeekMath~Corpus for training produces significant gains in mathematical reasoning for both English and Chinese tasks. Conversely, most existing datasets, which are heavily skewed toward English, either marginally enhance or potentially impede Chinese-language mathematical performance.
    \item \textbf{Large-scale}:
    DeepSeekMath~Corpus exceeds previous mathematical datasets in terms of volume by a considerable margin. Figure \ref{fig:corpora_comparisons} shows that DeepSeek-LLM 1.3B, when trained with DeepSeekMath~Corpus, benefits from a much sharper improvement trajectory and continues to advance as training proceeds. In comparison, baseline corpora—being significantly smaller—undergo repeated epochs during training, which rapidly leads model performance to stagnate.
\end{itemize}

\subsection{Training and Evaluating DeepSeekMath-Base 7B}

This section details DeepSeekMath-Base 7B, a foundational model characterized by robust reasoning, particularly within mathematical domains. The initialization of this model leverages DeepSeek-Coder-Base-v1.5 7B \citep{deepseek-coder}, after which it undergoes training on a corpus comprising 500B tokens. Data sources are apportioned as follows: 56\% derived from the DeepSeekMath Corpus, 4\% from AlgebraicStack, 10\% sourced from arXiv, 20\% representing code from Github, and the final 10\% consisting of natural language examples in both English and Chinese from Common Crawl. The training protocol largely mirrors that in Section \ref{sec:quality-policy}; however, the learning rate peak is adjusted to 4.2e-4, and the batch size is set at 10M tokens.

To thoroughly gauge the mathematical proficiency of DeepSeekMath-Base 7B, we systematically examine the model’s capacity for generating fully self-contained mathematical solutions, utilizing problem-solving tools, and carrying out formal theorem proving. We additionally assess general attributes of the base model, reporting on its natural language comprehension, logical reasoning, and programming aptitudes.

\paragraph{Mathematical Problem Solving with Step-by-Step Reasoning}

DeepSeekMath-Base is evaluated on its ability to resolve mathematical questions using few-shot chain-of-thought prompting \citep{cot}. The assessment spans eight benchmarks in both English and Chinese, encompassing quantitative reasoning datasets such as GSM8K \citep{gsm8k}, MATH \citep{MATH}, and CMATH \citep{wei2023cmath}, alongside multiple-choice tests like MMLU-STEM \citep{mmlu} and Gaokao-MathQA \citep{agieval}. These benchmarks collectively span a wide spectrum of mathematics, from elementary level up to college-level subject matter.

As indicated in Table \ref{tab:base_math_cot}, DeepSeekMath-Base 7B achieves top results across all eight benchmarks when compared with open-source base models—including the highly utilized Mistral 7B \citep{mistral} and the more recent Llemma 34B \citep{llemma}, the latter of which has undergone mathematical training on Proof-Pile-2 \citep{llemma}. In particular, DeepSeekMath-Base attains a margin exceeding 10\% absolute over existing open-source base models on the challenging MATH benchmark, and also outperforms Minerva 540B \citep{minerva}, a much larger closed-source base model (77-fold in size) based on PaLM \citep{palm} with specialized mathematical pretraining.

\paragraph{Mathematical Problem Solving with Tool Use} For assessing program-assisted mathematical reasoning, we perform evaluations on GSM8K and MATH with a few-shot program-of-thought prompting strategy \citep{pot,pal}. In this setting, models are guided to solve each question by composing a Python script, with access to libraries like \textit{math} and \textit{sympy} to support advanced computations. The answer is determined by executing the generated code. Results in Table \ref{tab:base_math_pot_proof} indicate that DeepSeekMath-Base 7B achieves superior performance compared to the previous leading Llemma 34B.

\paragraph{Formal Mathematics}

Automating formal proof generation enhances both reliability and efficiency in mathematical proof verification, a subject of growing interest in recent scholarship. We test DeepSeekMath-Base 7B using the informal-to-formal proving benchmark from \citep{dsp_proof}, where the objective is to construct a formal proof given an informal description, its formal expression, and an informal reasoning path. Evaluations are conducted on miniF2F \citep{minif2f}, a testbed for Olympiad-level formal mathematics, by having the model produce formal proofs in Isabelle with few-shot prompts. Consistent with \cite{dsp_proof}, we prompt the model for proof sketches, while leveraging Sledgehammer \citep{sledgehammer} to complete the remaining details via automated proving tools. Table \ref{tab:base_math_pot_proof} demonstrates that DeepSeekMath-Base 7B delivers robust performance for the autoformalization of proofs.

\paragraph{Natural Language Understanding, Reasoning, and Code}

The capabilities of the models on natural language understanding are measured by MMLU \citep{mmlu}, on reasoning with BBH \citep{bbh}, and on programming tasks with HumanEval \citep{codex} and MBPP \citep{mbpp}. Table \ref{tab:understanding_reasoning_code} shows that DeepSeekMath-Base 7B achieves marked improvements over its predecessor, DeepSeek-Coder-Base-v1.5 \citep{deepseek-coder}, in MMLU and BBH, demonstrating the advantageous effects of mathematical pretraining for language comprehension and reasoning skills. Moreover, due to the incorporation of code-related tokens during continual training, DeepSeekMath-Base 7B preserves the coding performance seen in DeepSeek-Coder-Base-v1.5 on HumanEval and MBPP. In summary, DeepSeekMath-Base 7B notably outperforms the generalist Mistral 7B model \citep{mistral} across all three evaluation benchmarks concerning reasoning and coding proficiency.

\section{Supervised Fine-Tuning}

\subsection{SFT Data Curation}

We develop a comprehensive instruction-tuning dataset in mathematics, encompassing a range of problem statements and solutions in both English and Chinese. This dataset draws from multiple mathematical domains and includes problems of assorted difficulty. Each problem is coupled with a solution written in the chain-of-thought (CoT) \citep{cot}, program-of-thought (PoT) \citep{pot,pal}, or tool-integrated reasoning style \citep{tora}. The curated dataset consists of 776,000 training samples.

\begin{itemize}[topsep=0pt]
    \item \textbf{English mathematical datasets}: Our annotation process involves GSM8K and MATH problems, for which we provide tool-integrated answers. We further incorporate a subset from MathInstruct \citep{MathInstruct}, in addition to the training data from Lila-OOD \citep{lila}, featuring solutions formulated as CoT or PoT. The English component encompasses varied mathematics disciplines, such as algebra, probability, number theory, calculus, and geometry.
    \item \textbf{Chinese mathematical datasets}: We gather mathematical questions used in Chinese K-12 curricula, covering 76 distinct sub-domains like linear equations. Solutions are prepared in both CoT and tool-integrated formats for these items.
\end{itemize}

\subsection{Training and Evaluating DeepSeekMath-Instruct 7B}

Here, we detail the training regimen for DeepSeekMath-Instruct 7B, an instruction-tuned variant built upon DeepSeekMath-Base. Training samples are concatenated at random up to a context limit of 4,000 tokens. The model is optimized over 500 training steps using a batch size of 256, with a constant learning rate set to $5 \times 10^{-5}$.

To assess mathematical capability, we measure performance with and without tool integration across four quantitative reasoning benchmarks, covering both English and Chinese tasks. We compare our results against several leading contemporary models:

\begin{itemize}[topsep=0pt]
    \item \textbf{Closed-source models} comprise: (1) members of the GPT series, with GPT-4 \citep{gpt4} and GPT-4 Code Interpreter~\footnote{\url{https://openai.com/blog/chatgpt-plugins\#\#code-interpreter}} representing peak capabilities, (2) Gemini Ultra and Pro \citep{gemini}, (3) Inflection-2 \citep{inflection-2}, (4) Grok-1~\footnote{\url{https://x.ai/model-card}}, as well as recent entrants from Chinese technology firms such as (5) Baichuan-3~\footnote{\url{https://www.baichuan-ai.com}}, and (6) the newest release in the GLM series, GLM-4~\footnote{\url{https://open.bigmodel.cn/dev/api\#glm-4}} \citep{glm}.
\end{itemize}

These models are designed for broad applications and have generally undergone multiple alignment steps.

\item \textbf{Open-source models} include both general-purpose and math-enhanced variants. The general-purpose set comprises: (1) DeepSeek-LLM-Chat 67B \citep{deepseek-llm}, (2) Qwen 72B \citep{qwen}, (3) SeaLLM-v2 7B \citep{seallm}, and (4) ChatGLM3 6B \citep{chatglm3}. Among models tailored for mathematical capabilities: (5) InternLM2-Math 20B~\footnote{\url{https://github.com/InternLM/InternLM-Math}}, built on InternLM2 and refined through mathematical data and instruction tuning; (6) Math-Shepherd-Mistral 7B, which leverages PPO training \citep{schulman2017proximal} on Mistral 7B \citep{mistral} with rewards from process supervision; (7) the WizardMath models \citep{wizardmath}, which augment Mistral 7B and Llama-2 70B \citep{llama2} for mathematical tasks via evolve-instruct (a variant of instruction tuning using AI-generated instructions) and PPO, with training examples mainly drawn from GSM8K and MATH; (8) MetaMath 70B \citep{metamath}, obtained by fine-tuning Llama-2 70B with a GSM8K and MATH-enriched corpus; (9) ToRA 34B \cite{tora}, produced by fine-tuning CodeLlama 34B for mathematical reasoning that incorporates tool usage; and (10) MAmmoTH 70B \citep{MathInstruct}, which involves instruction-tuning Llama-2 70B on MathInstruct.

Table \ref{tab:sft_rl_math} demonstrates that, when tool use is prohibited during evaluation, DeepSeekMath-Instruct 7B achieves outstanding performance in step-by-step mathematical reasoning. Specifically, on the advanced MATH benchmark, it exceeds the results of all other open-source and most proprietary models (including Inflection-2 and Gemini Pro) by a margin of at least 9\% in absolute terms. This advantage persists even against much larger models (like Qwen 72B) and models with specialized reinforcement learning on mathematics (such as WizardMath-v1.1 7B). Although DeepSeekMath-Instruct matches the performance of proprietary Chinese models such as GLM-4 and Baichuan-3 on the MATH dataset, it does not reach the levels of GPT-4 and Gemini Ultra.

When natural language reasoning is combined with program-based tool usage for answering questions, DeepSeekMath-Instruct 7B attains nearly 60\% accuracy on MATH, surpassing the performance of all known open-source models. On additional benchmarks, the model is on par with DeepSeek-LLM-Chat 67B—the former state-of-the-art open-source model, which is tenfold larger in scale.

\section{Reinforcement Learning}

\subsection{Group Relative Policy Optimization} Reinforcement learning (RL) has demonstrated effectiveness for enhancing the mathematical reasoning capabilities of LLMs after the initial Supervised Fine-Tuning (SFT) stage \citep{wang2023math,wizardmath}. Here, we present Group Relative Policy Optimization (GRPO), a reinforcement learning approach designed to be both efficient and effective.

\subsubsection{From PPO to GRPO} Proximal Policy Optimization (PPO) \citep{schulman2017proximal}, an actor-critic RL method, is frequently applied in the reinforcement learning fine-tuning phase for LLMs \citep{ouyang2022training}. PPO aims to optimize LLMs by maximizing the following surrogate objective:

\begin{equation}
\footnotesize
    \mathcal{J}_{PPO}(\theta) = \mathbb{E}{[q \sim P(Q), o \sim \pi_{\theta_{old}}(O|q)]} \frac{1}{|o|} \sum_{t=1}^{|o|} \min \left[ \frac{\pi_\theta(o_{t} | q, o_{<t})}{\pi_{\theta_{old}}(o_{t} | q, o_{<t})} A_{t}, \text{clip} \left( \frac{\pi_\theta(o_{t} | q, o_{<t})}{\pi_{\theta_{old}}(o_{t} | q, o_{<t})}, 1 - \varepsilon, 1 + \varepsilon \

Here, $\pi_{\theta}$ and $\pi_{\theta_{old}}$ denote the current and previous policy models, respectively, while $q$ and $o$ refer to questions and outputs drawn from the question dataset and the prior policy $\pi_{\theta_{old}}$. The parameter $\varepsilon$ is a clipping hyper-parameter employed in PPO to ensure stable learning dynamics. The term $A_t$ stands for the advantage, estimated using Generalized Advantage Estimation (GAE) \citep{gae}, with dependence on the set of rewards $\{r_{\ge t}\}$ and a learned value function $V_{\psi}$. Therefore, PPO mandates the training of a value function alongside the policy, and to prevent excessive optimization against the reward model, it is standard practice to augment the reward at each token with a KL penalty relative to a reference model \citep{ouyang2022training}, resulting in

\begin{equation}
    r_{t} = r_\varphi(q, o_{\le t}) - \beta \log\frac{\pi_{\theta}(o_{t}|q, o_{<t})}{\pi_{ref}(o_{t}|q, o_{<t})},
\label{eq:PPO-reward}
\end{equation}

where $r_\varphi$ is the learned reward model, $\pi_{ref}$ is a reference policy (often the initial SFT model), and $\beta$ is a scalar scaling the KL regularization.

Training a value function in PPO typically requires an additional network with a parameter count similar to the policy, incurring significant computational and memory overheads. Further, the value function acts as a baseline during RL optimization to reduce variance in the estimated advantages. In the context of LLMs, it is common for the reward model to provide scores only for the final token, posing challenges for accurate per-token value function training. To circumvent these issues, as illustrated in Figure \ref{fig:grpo}, we introduce Group Relative Policy Optimization (GRPO), which eliminates the need for training a separate value function as in PPO. Instead, GRPO utilizes the average reward over a set of outputs generated in response to a single question as the baseline. Concretely, for each question $q$, GRPO generates a batch of outputs $\{o_1, o_2, \cdots, o_G\}$ using $\pi_{\theta_{old}}$, and the policy is improved by maximizing

\begin{equation}
\footnotesize
\begin{split}
    \mathcal{J}_{GRPO}(\theta) &= \mathbb{E}{[q \sim P(Q), \{o_i\}_{i=1}^G \sim \pi_{\theta_{old}}(O|q)]}  \\
    & \frac{1}{G}\sum_{i=1}^G\frac{1}{|o_i|} \sum_{t=1}^{|o_i|} \left\{ \min \left[ \frac{\pi_\theta(o_{i,t} | q, o_{i,<t})}{\pi_{\theta_{old}}(o_{i,t} | q, o_{i,<t})} \hat{A}_{i,t}, \text{clip} \left( \frac{\pi_\theta(o_{i,t} | q, o_{i,<t})}{\pi_{\theta_{old}}(o_{i,t} | q, o_{i,<t})}, 1 - \varepsilon, 1 + \varepsilon \right)  \hat{A}_{i,t} \right] - \beta \mathbb{D}_{KL}\left[\pi_{\theta} || \pi_{ref}\right]\right\} ,
\end{split}
\label{eq:GRPO-obj}
\end{equation}

with $\varepsilon$ and $\beta$ remaining as tuning parameters, and $\hat{A}_{i,t}$ representing the group-relative advantage, which is computed using only the comparative rewards of outputs within the same batch—a procedure described in detail later. This method for estimating advantages is inherently compatible with reward models, which are typically trained on comparisons among responses to identical prompts. Unlike the reward structure in PPO, where KL regularization is embedded within the reward function, GRPO directly penalizes KL divergence between the current policy and a reference in the loss, simplifying the evaluation of $\hat{A}_{i,t}$. Furthermore, rather than employing the KL penalty defined in (\ref{eq:PPO-reward}), GRPO uses an unbiased estimator for the KL divergence as suggested by \citep{kl_approx}:

\begin{equation}
\small
    \mathbb{D}_{KL}\left[\pi_{\theta

\begin{algorithm}[t]
  \small
  \caption{Iterative Group Relative Policy Optimization}
  \textbf{Input} initial policy model $\pi_{\theta_{\text{init}}}$; reward models $r_\varphi$; task prompts $\mathcal{D}$; 
  hyperparameters $\varepsilon$, $\beta$, $\mu$
  \begin{algorithmic}[1]
    \State Initialize the policy model: $\pi_\theta \leftarrow \pi_{\theta_{\text{init}}}$
    \For{iteration = 1, \dots, I}
       \State Set the reference model: $\pi_{ref} \leftarrow \pi_{\theta}$
      \For{step = 1, \dots, M}
      \State Draw a minibatch $\mathcal{D}_b$ from $\mathcal{D}$
      \State Store the current policy: $\pi_{\theta_{old}} \leftarrow \pi_{\theta}$ 
      \State For each query $q \in \mathcal{D}_b$, generate $G$ responses $\{o_i\}_{i=1}^G \sim \pi_{\theta_{old}} (\cdot \mid q)$
      \State Evaluate the outputs $o_i$ using the reward model $r_{\varphi}$ to obtain $\{r_i\}_{i=1}^G$
      \State Use group relative advantage estimation to calculate $\hat{A}_{i,t}$ at the $t$-th position of $o_i$
      \For{GRPO iteration = 1, \dots, $\mu$}
        \State Update $\pi_{\theta}$ by optimizing the GRPO objective (Equation \ref{eq:GC-GRPO})
      \EndFor
    \EndFor \State Enhance $r_\varphi$ via ongoing training and a replay buffer. 
    \EndFor 
  \end{algorithmic}
  \textbf{Output} $\pi_\theta$
  \label{alg:iter-grpo}
\end{algorithm}

\subsubsection{Outcome Supervision RL with GRPO} Consider a question $q$ and a collection of $G$ outputs $\{o_1, o_2, \cdots, o_G\}$ generated from $\pi_{\theta_{old}}$. Each output is assessed with the reward model, producing $G$ corresponding rewards $\mathbf{r} = \{r_1, r_2, \cdots, r_G\}$. These reward values are standardized within the group by subtracting the group mean and dividing by the group standard deviation. In outcome supervision, the final normalized reward for each output $o_i$ serves as the reward for every token in the output. Formally, the advantage for any token $t$ in $o_i$ is set as $\hat{A}_{i, t} = \widetilde{r}_i = \frac{r_i- {\rm mean}(\mathbf{r})}{{\rm std}(\mathbf{r})}$. The policy is subsequently refined by maximizing the objective function in equation (\ref{eq:GRPO-obj}).

\subsubsection{Process Supervision RL with GRPO} Whereas outcome supervision delivers only terminal rewards for whole responses, it may be inadequate for guiding policies in elaborate mathematical problems. Motivated by \cite{wang2023math}, we also investigate process supervision, which assigns feedback after each intermediate reasoning step. For a query $q$ and $G$ sampled outputs $\{o_1, o_2, \cdots, o_G\}$, a process reward model computes step-wise rewards, resulting in the collection $\mathbf{R} = \{ \{r_1^{index(1)},\cdots,r_1^{index(K_1)}\}, \cdots,  \{r_G^{index(1)},\cdots,r_G^{index(K_G)}\} \}$, with $index(j)$ denoting the position of the $j$-th step's concluding token and $K_i$ representing the total steps for output $o_i$. These scores are then normalized by subtracting the overall mean and dividing by the standard deviation, as in $\widetilde{r}_i^{index(j)} = \frac{r_i^{index(j)} - {\rm mean(\mathbf{R})}}{{\rm std(\mathbf{R})}}$. Under this supervision scheme, the advantage at each token position $t$ is the aggregate of the normalized rewards for all subsequent steps, i.e., $\hat{A}_{i, t} = \sum_{index(j) \ge t} \widetilde{r}_i^{index(j)}$. The policy is then updated via

\subsubsection{Iterative RL with GRPO} As reinforcement learning progresses, previously trained reward models may become inadequate for effectively supervising the evolving policy model. To address this limitation, we investigate iterative reinforcement learning using GRPO. As outlined in Algorithm \ref{alg:iter-grpo}, this method involves repeatedly generating new datasets for reward model training by sampling outputs from the updated policy model. We continue training the original reward model via a replay strategy that maintains 10\% of the past data in the buffer. Subsequently, we designate the policy model as the reference and further optimize it using the refreshed reward model.

\subsection{Training and Evaluating DeepSeekMath-RL}

Our reinforcement learning experiments utilize DeepSeekMath-Instruct 7B as the foundation. The RL training dataset consists of approximately 144K questions in chain-of-thought format, covering problems from GSM8K and MATH sourced from the SFT corpus. To specifically assess RL’s effects on benchmarks lacking coverage during RL, we exclude all other SFT questions from training. The construction of the reward model’s training data follows the methodology proposed by \citep{wang2023math}. The initial reward model is trained using DeepSeekMath-Base 7B and a learning rate of 2e-5. For the GRPO process, the policy model’s learning rate is set to 1e-6, with a KL penalty coefficient of 0.04. For each sample, we generate $64$ responses. Maximum output length and training batch size are both set to 1024. The policy model receives only a single update after each exploration cycle. We benchmark DeepSeekMath-RL 7B using the same evaluation suite as DeepSeekMath-Instruct 7B. Here, tasks from GSM8K and MATH that involve chain-of-thought solutions are classified as in-domain, while all remaining benchmarks are treated as out-of-domain evaluations.

Table \ref{tab:sft_rl_math} presents a comparison between open- and closed-source models employing both chain-of-thought and tool-augmented reasoning approaches across English and Chinese benchmarks. Our observations are as follows: 1) When applying chain-of-thought reasoning, DeepSeekMath-RL 7B achieves accuracies of 88.2\% on GSM8K and 51.7\% on MATH. This level of performance exceeds that of any open-source model within the 7B–70B parameter range and outperforms most closed-source counterparts. 2) Importantly, DeepSeekMath-RL 7B was exclusively trained using instruction tuning data in the chain-of-thought format from GSM8K and MATH, with its initial weights sourced from DeepSeekMath-Instruct 7B. Even though the model was exposed to a narrowly focused dataset during training, it still yields superior results relative to DeepSeekMath-Instruct 7B on all evaluation metrics, thereby highlighting the advantages conferred by reinforcement learning.

\section{Discussion}
In this section, we elaborate on our insights derived from pre-training and RL experiments.

\subsection{Insights from Pre-Training}

We begin by detailing our observations regarding pre-training. Unless explicitly mentioned, all results discussed adhere to the experimental protocols specified in Section \ref{sec:quality-policy}. Notably, all references to the DeepSeekMath Corpus pertain to the 89B-token dataset produced in the second round of data preparation.

\subsubsection{The Role of Code Training in Enhancing Mathematical Reasoning}

A widely held but empirically underexplored assumption posits that training on code enhances a model’s reasoning capabilities. Our work aims to contribute partial empirical evidence to this claim, specifically within the realm of mathematical reasoning: incorporating code into training regimes significantly improves models’ performance on mathematical reasoning tasks, both with and without the use of external tools.

To investigate the effects of code pre-training on mathematical reasoning abilities, we evaluated models using both two-stage and one-stage training protocols:

\textbf{Two-Stage Training}

\begin{itemize}[topsep=0pt]
    \item \textbf{Code Training for 400B Tokens $\rightarrow$ Math Training for 150B Tokens}: DeepSeek-LLM 1.3B undergoes a training phase on 400B code tokens, followed by an additional stage involving 150B math tokens;
    \item \textbf{General Training for 400B Tokens $\rightarrow$ Math Training for 150B Tokens}: As a baseline comparison, we replace the initial code token training with 400B tokens from a broad, general-purpose corpus compiled by DeepSeek-AI, to evaluate whether code-specific data yields greater improvements in mathematical reasoning than general language data.
\end{itemize}

\textbf{One-Stage Training}

\begin{itemize}[topsep=0pt]
    \item \textbf{Math Training for 150B Tokens}: Here, DeepSeek-LLM 1.3B is trained exclusively on 150B math tokens;
    \item \textbf{Training on a mixture of 400B Code Tokens and 150B Math Tokens}: Since sequential math training after code training leads to diminished coding proficiency, we examine whether a single-stage mixed-token approach, in which math and code tokens are combined during training, both strengthens mathematical reasoning and reduces catastrophic forgetting.
\end{itemize}

\paragraph{Results}
Tables \ref{tab:code-math} and \ref{tab:code-math-general-eval} report downstream performance outcomes for the varied training regimens.

In both two-stage and mixed one-stage regimes, pretraining on code tokens is shown to enhance program-augmented mathematical reasoning. Table \ref{tab:code-math} illustrates that pretraining solely on code already provides substantial gains in tackling GSM8K and MATH datasets with Python-based solutions; subsequent dedicated math token training in the second phase results in further performance advances. Notably, the one-stage mixed-token setup ameliorates the problem of catastrophic forgetting observed in the two-stage approach, while also improving both coding ability (Table \ref{tab:code-math-general-eval}) and program-aided mathematical reasoning (Table \ref{tab:code-math}).

Moreover, code token pretraining benefits mathematical reasoning tasks even when no tool assistance is used. In the two-stage framework, code training alone confers moderate improvements, while also facilitating more effective subsequent math training—culminating in optimal performance. However, a one-stage mixture of code and math tokens results in decreased performance for pure mathematical reasoning tasks without tool use. We hypothesize that the limited capacity of the DeepSeek-LLM 1.3B model may impede its ability to jointly internalize the demands of both code and mathematics within a single-stage, blended training regimen.

\subsubsection{ArXiv Papers Appear Limited in Enhancing Mathematical Reasoning}

Inclusion of arXiv papers as a core element in math pre-training datasets is a common practice \citep{minerva,gpt-f,llemma,mathpile}. Nonetheless, there is a paucity of granular investigations into their specific influence on mathematical reasoning abilities. Somewhat unexpectedly, our experimental findings suggest that arXiv papers contribute little to boosting mathematical reasoning. Our study considers models of various scales, namely DeepSeek-LLM 1.3B and DeepSeek-Coder-Base-v1.5 7B \citep{deepseek-coder}, and utilizes arXiv corpora processed through different pipelines:

\begin{itemize}[topsep=0pt]
    \item \textbf{MathPile} \citep{mathpile}: This dataset encompasses 8.9B tokens curated using a series of heuristic rules for data cleaning and filtering, with more than 85\% sourced from scientific arXiv publications;
    \item \textbf{ArXiv-RedPajama} \citep{redpajama}: Consists of the full set of arXiv LaTeX documents, stripped of preambles, comments, macros, and bibliographies, amounting to a total of 28.0B tokens.
\end{itemize}

For our evaluation, DeepSeek-LLM 1.3B is trained for 150B tokens and DeepSeek-Coder-Base-v1.5 7B is trained for 40B tokens on each of the respective arXiv datasets. Our results indicate that arXiv-based corpora alone do not confer measurable advantages—and sometimes even lead to degraded performance—across an array of mathematical reasoning benchmarks of varying difficulty assessed in this work. The evaluation suite includes tasks in quantitative reasoning such as GSM8K and MATH (Table \ref{tab:arxiv-cot}), STEM-related multiple-choice tests like MMLU-STEM (Table \ref{tab:arxiv-cot}), as well as formal mathematics challenges including miniF2F (Table \ref{tab:arxiv_atp}).

Nevertheless, this observation is provisional and must be interpreted with caution. Our analysis has not encompassed:

\begin{itemize}[topsep=0pt]
    \item The potential role of arXiv tokens for mathematical tasks outside the present scope, for example, converting formalized theorems or proofs into their informal representations;
    \item The influence of incorporating arXiv tokens alongside other categories of training data;
    \item Whether the advantages of arXiv documents become apparent at larger parameter counts.
\end{itemize}

Therefore, these aspects merit additional inquiry, which we defer to subsequent research.

\subsection{Insights of Reinforcement Learning}
\subsubsection{Towards to a Unified Paradigm} Here, we introduce a comprehensive framework to examine multiple training regimes, including SFT, RFT, DPO, PPO, and GRPO, and present experiments designed to analyze components of this framework. Broadly, the gradient of a training method with respect to the parameter $\theta$ may be formulated as follows:

\begin{equation}
    \nabla_{\theta}\mathcal{J}_{\textcolor{red}{\mathcal{A}}}(\theta) = \mathbb{E}[\underbrace{(q,o) \sim \textcolor{red}{\mathcal{D}}}_{Data \ Source}]\left( \frac{1}{|o|} \sum_{t=1}^{|o|}  \underbrace{GC_{{\mathcal{A}}}(q, o, t, \textcolor{red}{\pi_{{rf}}})}_{Gradient \ Coefficient}  \nabla_{\theta}\log \pi_{\theta}(o_t | q, o_{<t})\right).
\label{eq:objective}
\end{equation}

This framework comprises three essential elements: 1) the \textit{Data Source} $\mathcal{D}$, responsible for specifying the set of training examples; 2) the \textit{Reward Function} $\pi_{{rf}}$, which delivers the feedback signal used for learning; and 3) the \textit{Algorithm} $\mathcal{A}$, which transforms both the data and the reward into a gradient coefficient $GC$ that modulates the strength and direction of learning signals. Leveraging this general structure, we examine several representative algorithms:

\begin{itemize}
    \item \textbf{Supervised Fine-tuning (SFT)}: This approach adapts a pretrained model through additional training on SFT datasets manually curated by humans.
    \item \textbf{Rejection Sampling Fine-tuning (RFT)}: Building upon the SFT model, RFT refines it further by training on model outputs that have been filtered, using responses to SFT-sourced queries, and selecting samples based on answer accuracy.
    \item \textbf{Direct Preference Optimization (DPO)}: DPO improves the SFT model by training with augmented outputs obtained from the SFT model, leveraging a pairwise DPO objective to guide optimization.
    \item \textbf{Online Rejection Sampling Fine-tuning (Online RFT)}: In contrast to standard RFT, Online RFT initializes the policy model from SFT and continually updates it by fine-tuning on outputs sampled in real time from the evolving policy.
    \item \textbf{PPO/GRPO}: These methods use the SFT model as the starting point for the policy and employ reinforcement learning by generating samples with the current policy for further optimization.
\end{itemize}

A summary of the components included in these approaches can be found in Table \ref{tab:data_policy}. For a comprehensive account of the derivation process, consult Appendix \ref{app:analysis-rl}.

\paragraph{Insights on Data Source} We classify data sources into two main types: online sampling and offline sampling. Online sampling refers to data collected from the ongoing exploration of the currently trained policy model, while offline sampling involves using data generated from the initial SFT model’s sampling outcomes. RFT and DPO are categorized as utilizing offline sampling, whereas Online RFT and GRPO employ online sampling strategies.

Referencing Figure \ref{fig:rl-analysis}, our findings indicate that Online RFT achieves markedly superior results compared to RFT across two benchmark tasks. In the early training phases, both Online RFT and RFT demonstrate similar performance, yet as training continues, Online RFT establishes a significant lead, underscoring the benefits of online sampling. This outcome aligns with expectations: at the outset, the actor and SFT model remain largely similar, resulting in sampled data that differ minimally. As training advances, however, data produced by the actor diverge more from those of the SFT model, thereby making real-time sampling increasingly beneficial.

\paragraph{Insights on Gradient Coefficient} In these algorithms, the reward signal is transformed into a gradient coefficient used to update model parameters. Our experiments distinguish between two types of reward functions: ‘Rule’ and ‘Model’. The ‘Rule’ approach evaluates response quality according to answer correctness, while the ‘Model’ approach employs a reward model trained to assign scores to responses based on labels derived from rule-based assessments. A crucial distinction, highlighted in Equations \ref{eq:GC-RFT} and \ref{eq:GC-GRPO}, is that GRPO dynamically tunes the gradient coefficient as a function of the reward model’s score, facilitating nuanced reinforcement and penalization of varying response magnitudes. In contrast, Online RFT does not incorporate such flexibility; it applies a uniform positive update to all correctly answered responses and does not impose penalties for incorrect answers.

As illustrated in Figure \ref{fig:rl-analysis}, GRPO outperforms online RFT, which underscores the effectiveness of adjusting both positive and negative gradient coefficients. Moreover, GRPO+PS yields better results than GRPO+OS, demonstrating the advantage of leveraging detailed, step-aware gradient coefficients. Additionally, our investigation into iterative RL—conducting two iterative cycles in our experiments—reveals that this approach yields marked performance gains, particularly notable during the first iteration, as depicted in Figure \ref{fig:iter-rl}.

\subsubsection{Why RL Works?} This study applies reinforcement learning to a select portion of instruction tuning data, resulting in notable improvements over the baseline instruction tuning model. To elucidate the mechanism behind the effectiveness of reinforcement learning, we assess Pass@K and Maj@K accuracies of both Instruct and RL models across two separate benchmarks. Results displayed in Figure \ref{fig:maj-pass} reveal that RL substantially boosts Maj@K scores, but has negligible effect on Pass@K. This suggests that RL mainly fortifies the robustness of the model's output distribution—improvements primarily emerge from elevating the frequency of correct answers in the TopK results rather than from enhancing intrinsic capabilities. Likewise, \citep{wang2023making} report a \textbf{misalignment problem} in reasoning tasks within SFT models, and demonstrate that SFT reasoning capabilities can be enhanced via various preference alignment strategies \citep{yuan2023rrhf,song2023preference,wang2023making}.

\subsubsection{How to Achieve More Effective RL?}
\label{sec:effec_rl}
Our findings indicate that reinforcement learning is highly effective for mathematical reasoning tasks. Furthermore, we propose a unified conceptual framework for understanding diverse representative training methodologies. Within this framework, all these methods can be regarded as forms of direct or simplified RL. Equation \ref{eq:objective} outlines three essential components: Data Source, Algorithm, and Reward Function. We discuss prospective future directions relating to each of these aspects.

\paragraph{Data Source} The data source constitutes the foundational input for all training approaches. For RL in particular, we define the data source as the collection of unlabeled questions paired with output samples drawn from the policy model. In this study, we limit our scope to the question set utilized in the instruction tuning phase and adopt a simple nucleus sampling strategy to generate responses. This limitation may explain why our RL pipeline primarily enhances Maj@K performance. Looking forward, we aim to investigate our RL approach using out-of-distribution prompts, paired with \textbf{advanced sampling (decoding) strategies}, such as those employing tree-search algorithms \citep{yao2023tree}. Additionally, the adoption of \textbf{efficient inference techniques} \citep{xia-etal-2023-speculative,leviathan2023fast,kwon2023efficient,xia2024unlocking}, which are crucial for maximizing policy model exploration efficiency, will be another focal point for future work.

\paragraph{Algorithms} Algorithms leverage both data and the feedback from the reward signal to compute gradient coefficients, which in turn adjust model parameters. In line with Equation \ref{eq:objective}, prevailing approaches typically \textbf{TRUST} the reward function's feedback to modulate the conditional probability assigned to particular tokens. Nevertheless, there is no guarantee that the reward signal is consistently dependable, particularly when facing highly intricate tasks. For instance, the PRM800K dataset \citep{lightman2023let}, despite its meticulous annotation by experienced annotators, still exhibits an error rate of approximately 20\% in its annotations\footnote{\url{https://github.com/openai/prm800k/issues/12\#issuecomment-1728491852}}. Consequently, we focus on reinforcement learning algorithms designed to withstand inaccuracies in reward signals. Our perspective is that \textbf{WEAK-TO-STRONG} \citep{burns2023weak} alignment strategies have the potential to revolutionize current learning algorithm paradigms.

\paragraph{Reward Function} The reward function constitutes the primary source of training feedback. Within RL frameworks, this function is frequently operationalized as a neural reward model. We identify three pivotal avenues for further development of reward models: 1) \textbf{Enhancing the generalization capabilities of the reward model.} To ensure that the reward model is effective across out-of-distribution inputs and sophisticated generation outputs is critical; failing this, RL may simply entrench the existing LLM output distributions rather than achieving substantive capability enhancement; 2) \textbf{Accurately representing the reward model's uncertainty.} Properly capturing uncertainty may serve as an essential interface connecting less robust reward models and weak-to-strong learning frameworks; 3) \textbf{Developing efficient, high-quality process reward models} capable of delivering granular feedback throughout complex reasoning tasks \citep{lightman2023let,wang2023math}.

\section{Conclusion, Limitation, and Future Work}
We introduce DeepSeekMath, an open-source model that surpasses all previous open-source approaches on the competitive MATH benchmark and attains performance levels close to those of proprietary systems. DeepSeekMath~builds upon DeepSeek-Coder-v1.5 7B, undergoing sustained training for 500B tokens, including a notable subset of 120B math-specific tokens derived from Common Crawl. Through a comprehensive ablation analysis, we identify that web pages provide a rich resource for mathematical data, while the anticipated benefits from arXiv are less significant than anticipated. Furthermore, we propose Group Relative Policy Optimization (GRPO), an adaptation of Proximal Policy Optimization (PPO), which significantly enhances mathematical reasoning ability while reducing memory usage. Experimental findings confirm the efficacy of GRPO, even when DeepSeekMath-Instruct 7B attains strong benchmark scores. Additionally, we propose a unified framework to contextualize various techniques, and outline several future avenues for advancing reinforcement learning effectiveness.

Despite DeepSeekMath's~notable achievements in quantitative reasoning, its performance remains comparatively limited in geometry and theorem-proving tasks when measured against closed-source models. In particular, our preliminary experiments reveal the model's difficulties with problems involving triangles and ellipses, potentially reflecting a bias in data selection during pre-training and fine-tuning stages. Furthermore, constrained by its parameter size, DeepSeekMath~does not match GPT-4 in few-shot learning capabilities: GPT-4 demonstrates performance gains with few-shot examples, while DeepSeekMath~delivers similar results in both zero-shot and few-shot settings. Going forward, our work will focus on enhancing the engineered data selection process to yield a superior pre-training corpus. We also plan to further investigate the promising approaches highlighted in Section \ref{sec:effec_rl} to improve reinforcement learning strategies for large language models (LLMs). 

\end{document}